\documentclass[10pt,a4paper]{amsart}
\usepackage[margin=19mm]{geometry}
\usepackage{amsmath,amssymb,mathtools}
\usepackage[scr=rsfs,cal=euler]{mathalfa}
\usepackage{xfrac}
\usepackage[unicode,hidelinks,bookmarks=false]{hyperref}
\newcommand{\ie}{i.e.\ }
\newcommand{\cf}{cf.\ }
\newcommand{\resp}{respectively\ }
\newcommand{\Subsection}{\subsection}
\providecommand{\nobreakdash}{\nobreak\hbox{-}}
\setlength{\emergencystretch}{2em}
\multlinegap0pt
\usepackage[shortlabels]{enumitem}
\usepackage{amssymb}
\newtheorem{theorem}{Theorem}
\newtheorem{proposition}{Proposition}
\newcommand{\F}{\mathbb{F}}
\newcommand\Fq{{\mathbb{F}_q}}
\DeclareMathOperator{\Ima}{Im}
\title{A Generalization of S\'ark\"ozy's theorem in function fields}
\author{Pierre-Yves Bienvenu, Th\'ai Ho\`ang L\^e,, Gauree Wathodkar}
\date{Source: arXiv:2609.12499v1 (CC BY 4.0)}
\begin{document}
\maketitle

\noindent\emph{Source and licence.} A Generalization of S\'ark\"ozy's theorem in function fields. Version arXiv:2609.12499v1 (CC BY 4.0), distributed by arXiv under the Creative Commons Attribution 4.0 International licence, http://creativecommons.org/licenses/by/4.0/, as recorded in the arXiv licence metadata for this exact version and shown on the abstract page https://arxiv.org/abs/2609.12499v1. Full licence notice: \texttt{licenses/arxiv-2609.12499-CC-BY-4.0.txt}.\par\smallskip

\noindent\textit{Editorial scope.} Counted unit: the article's proposition thmB (source0.tex lines 197--209). The article's complete original proof of it is delivered with it (source0.tex lines 212--264, one \texttt{proof} environment of 53 lines). No mathematics is written, repaired or re-derived: the statement and the proof are the article's own lines, in the article's own notation, and no retained line is substituted or re-flowed.  Retained uncounted as the source-local context the proof needs: lines 167--174.  Nothing was omitted from any retained span, so the package carries no omission record.  No retained line contains a citation, so the wrapper carries no bibliography and no bibliography entry.  No retained line contains a figure environment, an image-inclusion command or drawing code, so no image is delivered and the package declares no asset dependency.  Editorial formatting only, in addition: an AMS \texttt{amsart} preamble that restates the article's own declarations and macros used by the retained lines, namely the environments \texttt{theorem}, \texttt{proposition} on separate counters, together with the standard packages those lines need.  The retained environments print in this document as Theorem 1, Proposition 1 in order of appearance, whereas the article prints the same items with the numbering of its own full document.  The complete native source of the article as distributed in the arXiv e-print for this exact version is bundled unchanged as \texttt{sources/210120/source0.tex}.
\begin{theorem}\label{thmA}
	Let $r, k, d$ and $\ell$ be integers satisfying $k>2r^2 (1-1/\ell^2)$. Suppose $c_1,\dots , c_k\in \Fq[t]$ have degree at most $d$ and satisfy $\sum_{i=1}^{k}c_i=0$. Let $F(x) \in \Fq[t][x]$ be a polynomial of degree $\ell$ with constant term zero. 
 Then there exist constants $c = c(q,r, k ,\ell) \in (0,1)$ and $c'=c'(q,r,k,d,F)>0$ such that the following holds: Any $A\subset G_n$ satisfying $|A|> c' \cdot q^{cn} $ must contain a nontrivial solution to 
	\begin{equation}\label{ee}
		c_1 a_1^r+\dots +c_k a_k^r=F(b)
	\end{equation}
	for some $b\in \Fq[t]$.
\end{theorem}

\begin{proposition}\label{thmB}
Let $m,n,n',d'$ and $d''$ be positive integers. Let $\Psi: (\mathbb{F}_q^n)^k \rightarrow \mathbb{F}_q^{n'}$ be a polynomial map with degree at most $d'$ and $\Phi:\mathbb{F}_q^m \rightarrow \mathbb{F}_q^{n'}$ be a polynomial map with degree at most $d''$.
Let $A \subset \F_q^n$. Assume that
    \begin{enumerate}[label=(\roman*)]
        \item $\Psi^{-1}(0)\cap A^k = \{ (a, ..., a) : a \in A \}$,
        \item $|\Phi^{-1}(0)|$ is coprime to $q$, and
        \item $\Psi(A^k) \cap \Ima(\Phi) = \{0\}$.
    \end{enumerate}
    Then,
	\begin{equation} \label{eq:prop}
		|A|\leq k \cdot \inf_{0<x<1}\frac{(1+x+\dots+x^{q-1})^n}{x^{\left(\frac{q-1}{k}\right)\left(n'-\frac{m}{d''}\right)d' }}.
	\end{equation}
\end{proposition}

\begin{proof}[Proof of Theorem \ref{thmA} assuming Proposition \ref{thmB}]
By the observation made in the introduction, we can assume that $F(x)$ has only one root $x=0$ in $\Fq[t]$, by replacing $F(x)$ with $F(t^Mx)$ for some sufficiently large $M$.

We identify $G_n$ with $\F_q^n$, by identifying each polynomial $g=g_0+g_1t+\dots+ g_{n-1}t^{n-1} \in G_n$ with the vector $(g_0,g_1, \dots , g_{n-1})\in \mathbb{F}_q^n$. Then $g^r \in G_{(n-1)r+1}$ corresponds to the vector $(g_0^r, rg_0^{r-1}g_1,\ldots,g_{n-1}^r) \in \mathbb{F}_q^{r(n-1)+1}$. Note that each component of $g^r$ is a polynomial of degree at most $r$ in $g_0,g_1, \dots , g_{n-1}$. 
	
Under this identification, the map 
$
	(a_1,\dots, a_k) \mapsto \sum ^k_{i=1} c_i a_i^r
$
induces a polynomial map 
\begin{align*}
\Psi: (\mathbb{F}_q^n)^k & \rightarrow \mathbb{F}_q^{n'} \\  %, i.e. $d'=r$.
	(a_1,\dots, a_k) & \mapsto (\psi_1 (a_1,\dots, a_k), \ldots, \psi_{n'} (a_1,\dots, a_k) )
\end{align*}	
where $n':=(n-1)r+d+1$ and $\psi_i (a_1,\dots, a_k)$ is the coefficient of $t^{i-1}$ in $\sum ^k_{i=1} c_i a_i^r$. Moreover, the degree of each $\psi_i$ is at most $d':=r$.
	
Suppose 
\[
F(x) = \alpha_1(t) x+ \alpha_2(t) x^2+\cdots +\alpha_{\ell}(t) x^\ell \in \Fq[t][x].
\]
Let $u$ be the maximum of $\deg_t \alpha_{i}(t)$ for $i=1, \ldots, \ell$ and $m = \lfloor {\frac{n'-u}{\ell}}\rfloor +1$. If $h= h_0+h_1t+\cdots+h_{m-1}t^{m-1} \in G_m$, then 
	\begin{align*}
		F(h_0+h_1t&+\dots+h_{m-1}t^{m-1})\\
		&= f_0(h_0,\dots, h_{m-1})+f_1(h_0,\dots, h_{m-1})t+\dots+f_{n'-1}(h_0,\dots, h_{m-1})t^{n'-1} \in G_{n'}
	\end{align*}
	where $f_i \in \Fq[x_1, \dots,x_m]$ are polynomials with degrees at most $\ell$ and $f_i(0)=0$, for $0 \leq i\leq n'-1$.
	
Consider the polynomial map 
\begin{align*}
\Phi:\mathbb{F}_q^m & \rightarrow \mathbb{F}_q^{n'}, \\
(h_0, \ldots, h_{m-1}) &\mapsto (f_0(h_0,\dots, h_{m-1}), \ldots, f_{n'-1}(h_0,\dots, h_{m-1})). 
\end{align*}
	
Then $\Phi$ has degree at most $d'':=\ell$.  Suppose now $A \subset G_n$ is a subset that does not contain a nontrivial solution to $c_1 a_1 ^r + \cdots + c_k a_k^r = F(b)$.
We apply Proposition \ref{thmB} to $\Psi, \Phi$ and the set $A$. The first and third condition are satisfied by the assumption on $A$. The second condition is satisfied by the assumption on the number of roots in $\Fq[t]$ of $F$. 

%Therefore, Proposition \ref{thmB} implies that
%\begin{equation*}
%		|A|\leq k \inf_{0<x<1}\frac{(1+x+\dots+x^{q-1})^n}{x^{(q-1)(n'-m/d'')d'k^{-1} }}.
%\end{equation*}
    
Substituting $n'= (n-1)r+d+1, d'=r, d''=\ell$ and $m = \lfloor {\frac{n'-u}{\ell}}\rfloor +1$, \eqref{eq:prop} implies that
\[
|A|\leq k \inf_{0<x<1}\frac{(1+x+\dots+x^{q-1})^n}{ x^{\frac{q-1}{k} \left( 1 - \frac{1}{\ell^2} \right) nr^2 + O_{q,r,k,d,F}(1) }}.
\]
Writing $E := (q-1)r^2(1-1/\ell^2)k^{-1}$, then we have $E <(q-1)/2$ since $k>2r^2(1-1/\ell^2)$.

Let $f(x) =(1+x+\dots+x^{q-1})x^{-E}$. Then $f'(1) = \sum_{i=0}^{q-1} (i-E) > 0$, so $f$ is increasing on some interval $(1-\epsilon,1)$. Therefore, $\inf_{0<x<1} f(x) < f(1) = q$. Let $x_0 \in (0,1)$ be such that $f(x_0) <q$, then
  \begin{equation*}
|A|  \leq k \cdot x_0^{O_{q,r,k,d,F}(1)} \cdot (f(x_0))^n \leq c' \cdot q^{cn} 
  \end{equation*}
for some constants $c = c(q,r,k,\ell) \in (0,1)$ and $c'(q,r,k,d,F) >0$, as desired.
\end{proof}

\end{document}
